% Vista de contenedores (C4 nivel 2) del código fuente del sistema en su versión del 24 de octubre de 2026,
% con los flujos de la subida y los kits (A), el pago (B) y los avisos (C). Estilos locales: no depende del
% preámbulo; solo usa las bibliotecas arrows.meta y shapes.geometric.
\definecolor{tkPerF}{HTML}{FFF0B8}\definecolor{tkPerB}{HTML}{9A7B00}
\definecolor{tkAppF}{HTML}{CFE2F7}\definecolor{tkAppB}{HTML}{1F4E86}
\definecolor{tkProcF}{HTML}{D3EDD0}\definecolor{tkProcB}{HTML}{2E6B2B}
\definecolor{tkExtF}{HTML}{FAD7B5}\definecolor{tkExtB}{HTML}{A8500F}
\definecolor{tkInfF}{HTML}{DEDEDE}\definecolor{tkInfB}{HTML}{4D4D4D}
\definecolor{tkBdF}{HTML}{CDEBEB}\definecolor{tkBdB}{HTML}{1C6E6E}
\definecolor{tkMqF}{HTML}{E2D6F4}\definecolor{tkMqB}{HTML}{5B2E99}
\begin{tikzpicture}[x=1cm, y=1cm, font=\scriptsize,
  per/.style={draw, thick, rounded corners=7pt, align=center, draw=tkPerB, fill=tkPerF, inner sep=3pt, minimum height=7mm},
  cont/.style={draw, rounded corners=2pt, align=center, draw=tkAppB, fill=tkAppF, inner sep=2pt, minimum height=8mm},
  proc/.style={draw, rounded corners=2pt, align=center, draw=tkProcB, fill=tkProcF, inner sep=2pt, minimum height=7mm},
  ext/.style={draw, rounded corners=2pt, align=center, draw=tkExtB, fill=tkExtF, inner sep=2pt, minimum height=7mm},
  infra/.style={draw, rounded corners=2pt, align=center, draw=tkInfB, fill=tkInfF, inner sep=2pt, minimum height=7mm},
  bd/.style={draw, cylinder, shape border rotate=90, aspect=0.22, align=center, fill=tkBdF,
    minimum width=1.25cm, draw=tkBdB, minimum height=8mm, inner sep=1pt},
  grupo/.style={draw=black!55, rounded corners=3pt, fill=black!4},
  http/.style={-{Stealth[length=1.6mm]}, semithick},
  evt/.style={-{Stealth[length=1.6mm]}, densely dashed, semithick},
  cola/.style={-{Stealth[length=1.6mm]}, semithick, draw=tkMqB},
  hueco/.style={draw=white, line width=2.8pt},
  paso/.style={circle, draw, fill=white, inner sep=0.5pt, font=\scriptsize\bfseries, minimum size=4.4mm},
  rot/.style={font=\scriptsize, fill=white, inner sep=1pt},
  col/.style={anchor=north west, text width=4.9cm, align=left, inner sep=0pt, font=\scriptsize}]
% Personas y aplicaciones
\node[per, text width=2.6cm] (acudiente) at (2.45,0) {\textbf{Acudiente}};
\node[per, text width=3.8cm] (equipo) at (8.6,0) {\textbf{Equipo de Tukit}\\operador y administrador};
\node[cont, text width=3.4cm] (web) at (2.45,-1.45) {\textbf{Aplicación web progresiva}\\Next.js, puerto 3000};
\node[ext, text width=1.85cm] (google_places) at (5.5,-1.45) {\textbf{Google Maps}\\sugerencias de dirección};
\node[cont, text width=3.4cm] (admin) at (8.6,-1.45) {\textbf{Panel de administración}\\Next.js, puerto 3001};
\node[ext, text width=2.5cm] (wompi) at (13.4,-1.45) {\textbf{Wompi}\\pasarela de pago};
\node[cont, minimum width=14.3cm, text width=14cm] (gateway) at (7.35,-2.85) {\textbf{Puerta de enlace}
  (Python, puerto 8010): reenvía cada petición al servicio que atiende el prefijo de su ruta};
% Servicios que llaman a otros: operación, inteligencia y comercio
\draw[grupo] (0.2,-3.65) rectangle (4.55,-6.25);
\draw[grupo] (4.85,-3.65) rectangle (10.15,-6.25);
\draw[grupo] (10.45,-3.65) rectangle (14.95,-6.25);
\node[cont, text width=2.3cm] (operacion) at (1.6,-4.3) {\textbf{Operación}\\puerto 8015};
\node[bd] (base_operacion) at (3.7,-4.3) {base de\\operación};
\node[proc, text width=1.7cm] (instantanea) at (1.2,-5.55) {instantánea de métricas (a mano)};
\node[proc, text width=1.7cm] (repartidor) at (3.5,-5.55) {repartidor de avisos};
\node[cont, text width=2.1cm] (inteligencia) at (6.4,-4.3) {\textbf{Inteligencia}\\puerto 8013};
\node[bd] (base_inteligencia) at (9.2,-4.3) {base y cola\\de trabajos};
\node[proc, text width=1.45cm] (sondeador_listas) at (5.75,-5.55) {sondeador\\de listas};
\node[proc, text width=1.45cm] (trabajador) at (7.5,-5.55) {trabajador\\durable};
\node[proc, text width=1.45cm] (publicador_inteligencia) at (9.25,-5.55) {publicador\\de eventos};
\node[cont, text width=2.4cm] (comercio) at (11.85,-4.3) {\textbf{Comercio}\\puerto 8014};
\node[bd] (base_comercio) at (13.95,-4.3) {base de\\comercio};
\node[proc, text width=1.6cm] (barredor) at (11.5,-5.55) {barredor de reservas};
\node[proc, text width=1.45cm] (publicador_comercio) at (13.75,-5.55) {publicador\\de eventos};
% Bus de llamadas a identidad y catálogo, que no llaman a otros servicios
\draw[semithick] (2.95,-6.7) -- (12.55,-6.7);
\foreach \x in {2.95,8.55,12.55}{\draw[semithick] (\x,-6.25) -- (\x,-6.7);}
\draw[semithick] (10.3,-3.25) -- (10.3,-6.7);
\node[rot, font=\tiny, align=left, anchor=south west, inner sep=0.5pt] at (10.6,-6.68) {llamadas HTTP a\\identidad y catálogo};
% RabbitMQ: intercambio de eventos entre las dos filas de servicios
\node[draw=tkMqB, fill=tkMqF, rounded corners=2pt, minimum width=10.35cm, minimum height=6.5mm,
  inner sep=0pt] (rabbitmq) at (9.125,-7.25) {};
\node[align=center, inner sep=0pt, font=\scriptsize] at (5.95,-7.25) {\textbf{RabbitMQ}\\[-1.5pt]intercambio tukit.eventos};
% Servicios proveedores
\draw[grupo] (2.9,-7.8) rectangle (6.75,-10.3);
\node[cont, text width=1.8cm] (identidad) at (3.95,-8.45) {\textbf{Identidad}\\puerto 8011};
\node[bd] (base_identidad) at (6.0,-8.45) {base de\\identidad};
\node[proc, text width=2.4cm] (sondeador_colegios) at (4.75,-9.65) {sondeador de colegios};
\draw[grupo] (9.05,-7.8) rectangle (13.1,-10.3);
\node[cont, text width=2.0cm] (catalogo) at (10.2,-8.45) {\textbf{Catálogo}\\puerto 8012};
\node[bd] (base_catalogo) at (12.3,-8.45) {base de\\catálogo};
\node[proc, text width=2.0cm, fill=tkProcF!45, densely dashed] (herramientas) at (10.2,-9.65) {herramientas de\\curaduría (a mano)};
\node[proc, text width=1.3cm] (publicador_catalogo) at (12.15,-9.65) {publicador\\de eventos};
% Sistemas externos
\node[ext, text width=1.9cm] (gemini) at (7.85,-8.45) {\textbf{Gemini}\\lectura y vectores};
\node[ext, text width=1.7cm] (tiendas) at (14.4,-9) {\textbf{Tiendas en línea}\\productos y precios};
\node[ext, text width=1.9cm] (correo) at (1.25,-8.45) {\textbf{Correo saliente}\\SMTP};
\node[ext, text width=1.9cm] (whatsapp) at (1.25,-9.55) {\textbf{WhatsApp}\\enlace wa.me};
% Fila de infraestructura
\node[infra, text width=4.6cm] (postgres) at (4.9,-11.05) {\textbf{PostgreSQL 16 con pgvector}\\un servidor con las
  cinco bases;\\cada una con su bandeja de salida};
\node[infra, text width=4.2cm, densely dotted, anchor=north] (redis) at (9.6,-10.45) {\textbf{Redis}, opcional\\caché de vectores y de lecturas\\y contador del límite de tasa\\de la puerta de enlace};
\node[infra, text width=3.3cm, densely dotted] (observabilidad) at (13.55,-11.05) {\textbf{Observabilidad opcional}\\OpenTelemetry y Jaeger};
% Personas, aplicaciones y pasarela
\draw[http] (acudiente) -- (web);
\draw[http] (equipo) -- (admin);
\draw[http] (web) -- (google_places);
\draw[http] (web.south) -- (web.south |- gateway.north);
\draw[http] (admin.south) -- (admin.south |- gateway.north);
\draw[http] ([xshift=2mm]web.north west) -- (0.88,0.65) -| node[rot, pos=0.25, above] {abre el widget de pago de Wompi} (wompi.north);
\draw[http] (wompi.south) -- (wompi.south |- gateway.north) node[rot, midway, right] {aviso firmado};
\draw[http] ([xshift=-2mm]comercio.north east) -- ++(0,0.45) -- (15.25,-3.45) |- (wompi.east);
\draw[http] (gateway.south -| operacion) -- (operacion.north);
\draw[http] (gateway.south -| inteligencia) -- (inteligencia.north);
\draw[http] (gateway.south -| comercio) -- (comercio.north);
% Llamadas entre servicios que llaman
\draw[http] ([xshift=-2mm]operacion.north east) -- ++(0,0.35) -| ([xshift=3mm]inteligencia.north west);
\draw[http] ([xshift=-3mm]inteligencia.north east) -- ++(0,0.4) -| ([xshift=3mm]comercio.north west);
% Llamadas a identidad y catálogo desde el bus; la del catálogo cruza RabbitMQ
\draw[http] (3.6,-6.7) -- (3.6,0 |- identidad.north);
\draw[hueco] (10.3,-6.9) -- (10.3,-7.6);
\draw[http] (10.3,-6.7) -- (10.3,0 |- catalogo.north);
% Bandejas de salida que lee cada publicador
\draw[evt] (base_inteligencia.south) -- (base_inteligencia.south |- publicador_inteligencia.north);
\draw[evt] (base_comercio.south) -- (base_comercio.south |- publicador_comercio.north);
\draw[evt] (base_catalogo.south) -- (base_catalogo.south |- publicador_catalogo.north);
% Eventos por RabbitMQ: de los publicadores al intercambio y del intercambio a la cola de cada consumidor
\foreach \x in {4.1,5.75,7.1,9.25}{\draw[hueco] (\x,-6.58) -- (\x,-6.82);}
\draw[cola] (publicador_inteligencia.south) -- (9.25,0 |- rabbitmq.north);
\draw[cola] (publicador_comercio.south) -- (13.75,0 |- rabbitmq.north);
\draw[cola] (11.5,0 |- publicador_catalogo.north) -- (11.5,0 |- rabbitmq.south);
\draw[cola] (4.1,0 |- rabbitmq.north) -- (4.1,0 |- repartidor.south);
\draw[cola] (5.75,0 |- rabbitmq.north) -- (sondeador_listas.south);
\draw[cola] (7.1,0 |- rabbitmq.north) -- (7.1,0 |- trabajador.south);
\draw[cola] (5.2,0 |- rabbitmq.south) -- (5.2,0 |- sondeador_colegios.north);
% Llamadas a sistemas externos
\draw[hueco] (7.85,-6.58) -- (7.85,-6.82);
\draw[hueco] (7.85,-6.9) -- (7.85,-7.6);
\draw[http] (7.85,0 |- trabajador.south) -- (gemini.north);
\draw[http] (catalogo.west) -- (gemini.east);
\draw[http] (13.1,-9) -- (tiendas.west);
\draw[http] (2.75,0 |- repartidor.south) |- (correo.east);
\draw[semithick] (2.75,0 |- repartidor.south) |- (whatsapp.east);
% Pasos de los flujos
\node[paso] at (2.15,-2.15) {A1};
\node[paso] at (7.85,-7.73) {A1};
\node[paso] at (10.3,-6.05) {A2};
\node[paso] at (8.3,-2.15) {A3};
\node[paso] at (5.75,-6.45) {A4};
\node[paso] at (2.75,-2.15) {B1};
\node[paso] at (11.6,0.65) {B2};
\node[paso] at (13.05,-2.15) {B3};
\node[paso] at (barredor.north west) {B4};
\node[paso] at (4.1,-6.45) {C1};
\node[paso] at (2.95,-6.45) {C2};
\node[paso] at (2.75,-9) {C3};
% Pasos de los flujos y notación, debajo de la figura
\node[col] (colA) at (0.2,-12.1) {\textbf{A. Subida de la lista y kits}\\[1pt]
  \textbf{A1} El acudiente sube la lista; inteligencia la lee con Gemini, guarda el trabajo y lo anuncia en su bandeja de salida.\\[1pt]
  \textbf{A2} Con la subida aceptada (202), el trabajador toma el trabajo de RabbitMQ o de su tabla, pide candidatos al catálogo y deja los kits por revisar.\\[1pt]
  \textbf{A3} El equipo revisa los kits en el panel y los publica; inteligencia deja el evento de kits publicados en su bandeja de salida.\\[1pt]
  \textbf{A4} Si la lista no estaba verificada, el sondeador de listas arma el kit al recibir el evento de lista verificada.};
\node[col] (colB) at (5.35,-12.1) {\textbf{B. Pago}\\[1pt]
  \textbf{B1} El acudiente elige un kit y escribe la dirección con sugerencias de Google Maps; inteligencia crea la orden en comercio, que reserva el inventario al precio vigente.\\[1pt]
  \textbf{B2} La aplicación web abre el widget de Wompi y el acudiente paga.\\[1pt]
  \textbf{B3} Wompi envía el aviso firmado o la página de resultado pide sincronizar el pago; comercio deja el evento de pedido pagado en su bandeja.\\[1pt]
  \textbf{B4} El barredor libera la reserva del pedido que no se paga en 24 horas y lo cancela.};
\node[col] (colC) at (10.5,-12.1) {\textbf{C. Avisos}\\[1pt]
  \textbf{C1} El repartidor de avisos recibe por RabbitMQ los eventos de kits publicados, pago, incidencia, despacho, reparto y entrega, y no repite el aviso de un evento ya enviado.\\[1pt]
  \textbf{C2} Pide a identidad el contacto del acudiente, sin guardarlo.\\[1pt]
  \textbf{C3} Si el acudiente tiene celular, arma un enlace de WhatsApp con el aviso; si no, se lo envía por correo.\\[1pt]
  El despacho lo registra el equipo en el panel, con el nombre y el celular del repartidor propio; las transportadoras con guía son una evolución prevista.};
\node[anchor=north west, text width=15.1cm, align=left, inner sep=0pt, font=\scriptsize]
  at ([yshift=-2mm]colA.south west |- current bounding box.south) {Línea continua: llamada HTTP. Línea discontinua:
  el publicador lee la bandeja de salida de su base. Línea violeta: evento que viaja por RabbitMQ hasta la cola de su
  consumidor (avisos, colegios, listas y trabajos), con reintentos y descartados; sin RabbitMQ, los publicadores no
  arrancan y cada consumidor lee la bandeja del productor. Línea punteada: componente opcional. Línea sin punta: enlace que el sistema
  arma y no envía. Recuadro gris: el servicio con sus procesos y su base.};
% Leyenda
\def\tkley#1#2#3#4{\node[#1, minimum width=7mm, minimum height=4.2mm, inner sep=0pt, text width=7mm] (l#2) at #3 {};
  \node[anchor=west, font=\scriptsize, inner sep=1pt] at ([xshift=1mm]l#2.east) {#4};}
\node[draw=black!40, rounded corners=2pt, anchor=north west, minimum width=15.05cm, minimum height=2.3cm,
  inner sep=0pt] (cajaley) at ([yshift=-1mm]colA.south west |- current bounding box.south) {};
\node[anchor=north west, font=\scriptsize\bfseries] at ([shift={(2mm,-1mm)}]cajaley.north west) {Leyenda};
\begin{scope}[shift={(cajaley.north west)}]
\tkley{per}{a}{(1.2,-0.75)}{Persona}
\tkley{cont}{b}{(3.9,-0.75)}{Aplicación o servicio}
\tkley{proc}{c}{(7.8,-0.75)}{Proceso en segundo plano}
\tkley{ext}{d}{(12.3,-0.75)}{Sistema externo}
\tkley{bd}{e}{(1.2,-1.35)}{Base de datos}
\tkley{infra}{f}{(3.9,-1.35)}{Infraestructura}
\node[draw=tkMqB, fill=tkMqF, rounded corners=2pt, minimum width=7mm, minimum height=4.2mm, inner sep=0pt] (lg) at (7.8,-1.35) {};
\node[anchor=west, font=\scriptsize, inner sep=1pt] at ([xshift=1mm]lg.east) {RabbitMQ};
\draw[http] (0.8,-1.95) -- +(0.9,0); \node[anchor=west, font=\scriptsize, inner sep=1pt] at (1.8,-1.95) {HTTP};
\draw[evt] (3.9,-1.95) -- +(0.9,0); \node[anchor=west, font=\scriptsize, inner sep=1pt] at (4.9,-1.95) {Evento de la bandeja};
\draw[cola] (8.4,-1.95) -- +(0.9,0); \node[anchor=west, font=\scriptsize, inner sep=1pt] at (9.4,-1.95) {Cola de RabbitMQ};
\end{scope}
\end{tikzpicture}
